\documentclass[11pt,letterpaper]{article}

% ── Packages ──────────────────────────────────────────────────────────────────
\usepackage[margin=1in]{geometry}
\usepackage{graphicx}
\usepackage{booktabs}
\usepackage{amsmath}
\usepackage{hyperref}
\usepackage{caption}
\usepackage{subcaption}
\usepackage{float}
\usepackage[numbers]{natbib}
\usepackage{authblk}

% ── Title ─────────────────────────────────────────────────────────────────────
\title{Deep Learning for Single-Cell RNA-seq Cell Type Classification:\\
       A 1D CNN Baseline, Pathway-Aware Transformer, and Graph Neural Network}

\author{Progress Report 2}
\affil{Deep Learning Final Project}
\date{\today}

\begin{document}
\maketitle

% ── Abstract ──────────────────────────────────────────────────────────────────
\begin{abstract}
Automated cell type annotation from single-cell RNA sequencing (scRNA-seq) data
is a critical bottleneck in modern transcriptomics. We present a three-way
comparison of deep learning architectures for this task using the Allen Brain
Atlas Human Primary Motor Cortex (M1) 10x Chromium dataset (76,533 cells,
50,281 genes, 17 cell types). First, a 1D convolutional neural network (CNN)
baseline operating on highly variable genes (HVGs) achieves 66.14\% overall
accuracy and 92.51\% balanced accuracy. Second, a pathway-masked Transformer
(TOSICA-inspired) operating on 300 Reactome pathway tokens achieves 66.56\%
overall accuracy and 92.65\% balanced accuracy, with interpretable per-pathway
attention scores and a built-in unknown-cell detection mechanism that flags
95.2\% of SMART-seq transfer cells as low-confidence. Third, a GraphSAGE
graph neural network (GNN) operating on a cell-similarity $k$-NN graph achieves
63.17\% overall accuracy and 90.05\% balanced accuracy, with improved handling
of the problematic L2/3~IT class (F1 0.247 vs.\ 0.011 for the CNN) and 88.56\%
accuracy on inductive SMART-seq transfer. Analysis of failure modes across all
three architectures reveals that the L2/3~IT/Micro-PVM misclassification
pathology is a data-level phenomenon driven by class imbalance, not an artifact
of any single model family.
\end{abstract}

% ══════════════════════════════════════════════════════════════════════════════
\section{Introduction}

Single-cell RNA sequencing enables transcriptomic profiling at cellular
resolution, but the resulting datasets routinely contain tens of thousands of
cells measured across tens of thousands of genes. Manual cell type annotation
via marker-gene inspection after unsupervised clustering remains the dominant
workflow, yet it is subjective, labor-intensive, and does not scale to atlas-level
datasets. Supervised deep learning offers an alternative: train a classifier on
a well-annotated reference atlas and transfer predictions to new data.

We investigate this paradigm using the Allen Institute's Human M1 10x Chromium
atlas as the reference. Our primary contribution is a thorough 1D~CNN baseline
with honest evaluation, including per-class F1 analysis that exposes how
headline accuracy can mask catastrophic failures on individual classes. We then
evaluate a TOSICA-inspired Transformer architecture that addresses several
shortcomings of the CNN through pathway-level token embeddings and
attention-based interpretability. Finally, we develop a GraphSAGE-based GNN
that explicitly models inter-cell relationships through a cosine-similarity
$k$-NN graph, providing a third perspective on the classification problem and
a natural inductive transfer mechanism.

% ══════════════════════════════════════════════════════════════════════════════
\section{Data}

\subsection{Primary Dataset}
The Allen Brain Atlas Human Primary Motor Cortex (M1) 10x Chromium dataset
comprises 76,533 labelled cells and 50,281 genes. After filtering classes with
fewer than 100 cells, 76,362 cells spanning 17 subclass-level cell types were
retained. The data exhibit severe class imbalance (Figure~\ref{fig:dist}):
L2/3~IT neurons alone account for 19,385 training cells (31.7\%), while
Micro-PVM contributes only 86 (0.1\%).

\subsection{Secondary Dataset}
We additionally evaluate cross-technology transfer on the Allen Institute's
Human Multiple Cortical Areas SMART-seq dataset. This dataset uses plate-based
full-length sequencing rather than droplet-based 3$'$ counting, introducing
substantial technical batch effects.

\subsection{Practical Data Challenges}
The expression matrix (76K $\times$ 50K, ${\sim}$30\,GB as CSV) could not be
loaded into memory with standard \texttt{pandas.read\_csv}, even on Google
Colab's high-RAM runtime. We briefly considered provisioning a Google Cloud VM
with sufficient RAM but decided against the cost and complexity. Instead, we
adopted Polars lazy-frame streaming (\texttt{pl.scan\_csv} with
\texttt{engine="streaming"}), which processes the file in bounded memory by
evaluating filter, projection, and aggregation operations in a single pass
without materializing the full DataFrame. An earlier attempt to download the
matrix via \texttt{curl~-o} terminated prematurely, producing a truncated file
that happened to yield artificially strong classification performance---a
cautionary lesson in validating data integrity before drawing conclusions.

% ══════════════════════════════════════════════════════════════════════════════
\section{Methods}

\subsection{Preprocessing}
Data were split 80/10/10 into train, validation, and test sets \emph{before}
any preprocessing to prevent data leakage. The preprocessing pipeline operated
in three streaming phases, each scanning the CSV without loading it fully into
memory:

\begin{enumerate}
    \item \textbf{Gene filter (Phase~1).} Using only validation-set cells,
    genes expressed in fewer than 3\% of cells were removed, retaining 12,402
    of 50,281 genes.
    \item \textbf{HVG selection (Phase~2).} On the same validation cells,
    library-size normalization ($\times 10^{4}$) and log$_{1p}$
    transformation were applied in-stream, and per-gene variance was computed.
    The top 1,000 genes by variance were selected as highly variable genes
    (HVGs).
    \item \textbf{Data extraction (Phase~3).} A final streaming pass over all
    76,362 cells extracted and normalized only the 1,000 HVG columns, then
    scattered rows into pre-allocated per-split NumPy arrays indexed by a
    sample-name lookup table.
\end{enumerate}

\noindent Finally, $z$-scoring via \texttt{StandardScaler} (fit on training
data only) completed the CNN and GNN pipelines. The Transformer pipeline used
library-size normalization and log$_{1p}$ only, without $z$-scoring, consistent
with the TOSICA design. Intermediate Polars DataFrames were explicitly deleted
and garbage-collected after each phase to minimize peak memory.

\subsection{1D CNN Architecture}
Gene expression vectors were treated as single-channel 1D signals of length
1,000. The feature extractor comprised four stages:

\begin{enumerate}
    \item \textbf{Conv1d}$(1 \to 64,\; k{=}15)$ with same-padding
    \item \textbf{ConvBlock}$(64 \to 128,\; k{=}15) \to$ MaxPool$(4)$
    \item \textbf{ConvBlock}$(128 \to 256,\; k{=}9) \to$ MaxPool$(4)$
    \item \textbf{ConvBlock}$(256 \to 128,\; k{=}7) \to$ MaxPool$(4)$
\end{enumerate}

\noindent Each ConvBlock applies Conv1d $\to$ BatchNorm $\to$ Dropout(0.1) $\to$
GELU\@. An adaptive average pool reduces the final feature map to a 128-d vector,
which feeds a classifier head: Linear$(128 \to 200) \to$ ReLU $\to$
BatchNorm $\to$ Linear$(200 \to 17)$. Total parameters: 678,769.

\subsection{Pathway-Masked Transformer Architecture}
Rather than selecting HVGs by variance, the TOSICA-inspired model constructs a
binary mask $\mathbf{M} \in \{0,1\}^{G \times P}$ from MSigDB Reactome gene
sets, where $M_{gp}=1$ iff gene $g$ belongs to pathway $p$. After intersecting
with the 10x gene universe, 8,822 genes across 300 pathways were retained. The
masked embedding layer computes per-pathway tokens as:

\[
\mathbf{e}_{b,p} = \sum_{g} x_{b,g}\,(\mathbf{W}_{:,g,p} \odot M_{g,p}) + \mathbf{b}_p
\]

\noindent A learnable CLS token is prepended, positional embeddings are added,
and the sequence of $P{+}1$ tokens is processed by a 2-layer Transformer
encoder (embedding dim 48, 4 heads, pre-norm). The CLS output feeds a
LayerNorm $\to$ Linear$(48 \to 96) \to$ GELU $\to$ Dropout $\to$ Linear$(96
\to 17)$ classifier. Cells with $\max(\mathbf{p}) < 0.95$ are flagged as
``Unknown.'' Total parameters: $\sim$150K.

\subsection{GraphSAGE GNN Architecture}
The GNN explicitly models inter-cell relationships. A cosine $k$-NN graph
($k{=}15$) was constructed over all 76,362 training cells on their 200
top-HVG features, yielding 1,858,284 symmetric edges. Node features are the
200-dimensional $z$-scored HVG vectors.

The architecture uses two ResidualSAGEBlocks followed by a final SAGEConv:

\[
\mathbf{h}^{(l+1)} = \text{GELU}(\text{LayerNorm}(\text{SAGEConv}(\mathbf{h}^{(l)}, \mathcal{E}))) + \mathbf{W}_\text{proj}\,\mathbf{h}^{(l)}
\]

\noindent The final hidden dimension is 256 (halved to 128 after the last conv).
A LayerNorm $\to$ Linear$(128 \to 17)$ head produces logits. For inductive
SMART-seq transfer, SMART-seq node features are stacked with training features,
new cells are connected to their $k$ nearest training neighbors, and a single
forward pass produces embeddings for all nodes simultaneously. Total parameters:
$\sim$290K.

\subsection{Training Protocol}
All models were trained on Google Colab using an NVIDIA T4 GPU with the
high-RAM runtime setting. Training used AdamW
($\text{lr}{=}3{\times}10^{-4}$, weight decay$\,{=}10^{-6}$), cosine
annealing, class-weighted cross-entropy with label smoothing of 0.1, gradient
clipping at 1.0, and mixed-precision (FP16) on CUDA\@. Early stopping with
patience 15 monitored validation loss. The GNN converged at epoch 33.

% ══════════════════════════════════════════════════════════════════════════════
\section{Results: 1D CNN Baseline}

\subsection{Exploratory Visualization}
Figure~\ref{fig:dist} shows the severe class imbalance in the training set.
PCA on the 1,000 HVGs (Figure~\ref{fig:pca}) captures only ${\sim}$57\%
cumulative variance in 50 components, with PC1 (14.6\%) and PC2 (8.6\%)
already separating major non-neuronal populations (Oligo, Pvalb) from
excitatory neurons. The UMAP embedding (Figure~\ref{fig:umap}) further
resolves all 17 classes into visually distinct clusters, confirming that
the HVG representation carries discriminative signal. A dendrogram-ordered
heatmap of the top 100 HVGs (Figure~\ref{fig:heatmap}) reveals
cell-type-specific expression blocks, with clear marker patterns for
non-neuronal types (Astro, OPC, Oligo) and more subtle distinctions among
excitatory layer subtypes.

\begin{figure}[H]
    \centering
    \includegraphics[width=\linewidth]{figures/cell-type-distibution.png}
    \caption{Training set cell type distribution. L2/3~IT dominates at 19,385
    cells; Micro-PVM has only 86.}
    \label{fig:dist}
\end{figure}

\begin{figure}[H]
    \centering
    \includegraphics[width=\linewidth]{figures/PCA.png}
    \caption{Left: PCA scree plot showing cumulative variance does not reach
    80\% within 50 components, reflecting the high dimensionality of
    scRNA-seq data. Right: PC1 vs.\ PC2 colored by cell type.}
    \label{fig:pca}
\end{figure}

\begin{figure}[H]
    \centering
    \includegraphics[width=0.65\linewidth]{figures/UMAP.png}
    \caption{UMAP of training set (1,000 cells subsampled) on 50 PCs. All
    17 cell types form visually separable clusters.}
    \label{fig:umap}
\end{figure}

\begin{figure}[H]
    \centering
    \includegraphics[width=\linewidth]{figures/den_heatmap-of-gene-expression.png}
    \caption{Dendrogram-clustered heatmap of the top 100 HVGs across 50 cells
    per type. Row color bars indicate cell type; columns are hierarchically
    clustered genes. Red/blue encodes $z$-scored expression.}
    \label{fig:heatmap}
\end{figure}

\subsection{Training Dynamics}
The loss and accuracy curves (Figure~\ref{fig:loss}) reveal notable instability.
Training accuracy spikes to 94.8\% by epoch~3, then collapses to ${\sim}$64\%
by epoch~6 and remains there. Validation accuracy oscillates between 39\% and
75\% in the first 10 epochs before stabilizing near 66\%. This behavior is
consistent with the model initially learning to classify many classes well, then
converging to a degenerate solution in which the largest class (L2/3~IT) is
systematically misclassified---an equilibrium the class-weighted loss and cosine
annealing cannot escape within 20 epochs.

\begin{figure}[H]
    \centering
    \includegraphics[width=\linewidth]{figures/loss_curve_for_one_cnn.png}
    \caption{Training and validation loss/accuracy curves over 20 epochs.
    The sharp training accuracy collapse after epoch~3 suggests the model
    shifts toward a degenerate prediction mode for the dominant class.}
    \label{fig:loss}
\end{figure}

\subsection{Test Set Evaluation}
The model achieves 66.14\% overall test accuracy, 92.51\% balanced accuracy,
and 0.8705 macro F1 (Table~\ref{tab:comparison}). The disparity between overall
accuracy and balanced accuracy is diagnostic: the model achieves near-perfect
recall ($>$0.94) on 15 of 17 classes but catastrophically fails on L2/3~IT
(recall 0.005, F1 0.011) and Micro-PVM (precision 0.004, F1 0.009).

The confusion matrix (Figure~\ref{fig:cm}) shows that nearly all L2/3~IT cells
(99\%) are misclassified as Micro-PVM, and conversely, Micro-PVM achieves
perfect recall but absorbs thousands of false positives---hence its precision of
0.004. This cross-contamination between the largest and smallest classes is the
primary driver of the low headline accuracy. All remaining classes achieve
F1\,$\geq$\,0.95 (Figure~\ref{fig:f1}).

\begin{figure}[H]
    \centering
    \includegraphics[width=0.85\linewidth]{{figures/66_acc_on_test_set_Confusion matrix}.png}
    \caption{Normalized confusion matrix on the test set. The L2/3~IT row shows
    near-total misclassification into Micro-PVM (column 10), the single failure
    mode responsible for the 66\% overall accuracy.}
    \label{fig:cm}
\end{figure}

\begin{figure}[H]
    \centering
    \includegraphics[width=\linewidth]{figures/f1-per-classgroup.png}
    \caption{Per-class F1 scores. The mean F1 of 0.870 is dragged down
    almost entirely by L2/3~IT (0.01) and Micro-PVM (0.01); all other classes
    exceed 0.95.}
    \label{fig:f1}
\end{figure}

\subsection{Cross-Technology Transfer}
The trained CNN was applied to the SMART-seq dataset by aligning the same 1,000
HVGs (zero-filling missing genes), applying identical normalization, and running
inference. Without explicit batch correction, the model must generalize across
fundamentally different capture technologies. ComBat correction was attempted
but library dependencies were unavailable in the runtime.

% ══════════════════════════════════════════════════════════════════════════════
\section{Results: Pathway-Masked Transformer}

\subsection{Training Dynamics}
Figure~\ref{fig:attn_loss} shows the Transformer training curves. The model
trains stably without the collapse seen in the CNN, reaching convergence within
the 100-epoch budget. The smooth loss trajectory reflects the pathway-level
aggregation, which acts as a built-in regularizer by preventing the model from
overfitting to individual gene noise.

\begin{figure}[H]
    \centering
    \includegraphics[width=\linewidth]{figures/attention_loss_and_accuracy.png}
    \caption{Transformer training and validation loss/accuracy curves.
    Training is notably more stable than the CNN, without the early-epoch
    accuracy collapse observed in the convolutional baseline.}
    \label{fig:attn_loss}
\end{figure}

\subsection{Test Set Evaluation}
The Transformer achieves 66.56\% overall test accuracy, 92.65\% balanced
accuracy, and 0.8672 macro F1 (Table~\ref{tab:comparison})---a marginal but
consistent improvement over the CNN on every aggregate metric. Critically, the
L2/3~IT/Micro-PVM misclassification pattern persists: L2/3~IT recall remains
near zero, confirming that this failure is a data-level phenomenon (extreme
class imbalance and feature overlap) rather than an architectural one.

\subsection{Attention-Based Interpretability}
Because the Transformer encodes cells as sequences of pathway tokens, the
CLS-to-pathway attention weights constitute a biologically meaningful cell
embedding. UMAP on these attention scores (Figure~\ref{fig:attn_umap}) produces
a visualization where proximity reflects shared pathway activity rather than
abstract statistical variance. Non-neuronal populations (Astro, Oligo, OPC)
form tight, well-separated clusters consistent with their distinct transcriptional
programs; excitatory layer subtypes (L4~IT, L5~IT, L6~IT) occupy adjacent
regions reflecting their shared cortical identity.

\begin{figure}[H]
    \centering
    \includegraphics[width=0.85\linewidth]{figures/umap_on_attention_score.png}
    \caption{UMAP of per-cell CLS-to-pathway attention scores. Each axis
    corresponds to a combination of Reactome pathway activities rather than
    a principal component of raw gene expression. Unknown SMART-seq cells
    (gray) cluster at the periphery, consistent with technical batch effects
    displacing them from the reference manifold.}
    \label{fig:attn_umap}
\end{figure}

\subsection{Cross-Technology Transfer and Unknown Detection}
When applied inductively to the SMART-seq dataset, the Transformer's
probability-threshold mechanism flags 95.2\% of cells as ``Unknown''
($\max(\mathbf{p}) < 0.95$). This high rejection rate reflects the substantial
technology-driven distributional shift between 10x droplet and SMART-seq
plate-based sequencing; the pathway tokenization alone is insufficient to bridge
this gap without explicit batch correction. The Unknown cells appear at the
periphery of the attention-score UMAP (gray points in
Figure~\ref{fig:attn_umap}), visually confirming their displacement from the
reference manifold.

% ══════════════════════════════════════════════════════════════════════════════
\section{Results: GraphSAGE Graph Neural Network}

\subsection{Graph Construction}
A cosine $k$-NN graph ($k{=}15$) was built over all 76,362 training cells on
their 200-HVG $z$-scored features, yielding 1,858,284 symmetric directed edges.
The graph captures local cell-similarity structure that neither the CNN nor the
Transformer explicitly models: neighboring nodes in the graph are cells with
similar expression profiles, and the GNN aggregates this neighborhood context
before making a classification decision.

\subsection{Training Dynamics}
Figure~\ref{fig:gnn_loss} shows the GNN training curves. The model converged
at epoch 33 via early stopping, substantially faster than the CNN (20 full
epochs) and Transformer (up to 100 epochs). The full-batch training
regime---processing all 76,362 nodes and 1.86M edges in a single forward
pass---is memory-intensive but avoids the mini-batch boundary artifacts that
can affect graph models trained with neighborhood sampling.

\begin{figure}[H]
    \centering
    \includegraphics[width=\linewidth]{figures/GGN_loss_and_accuracy.png}
    \caption{GNN training and validation loss/accuracy curves. Early stopping
    at epoch 33 indicates rapid convergence on the full-graph objective.}
    \label{fig:gnn_loss}
\end{figure}

\subsection{Test Set Evaluation}
The GNN achieves 63.17\% overall test accuracy, 90.05\% balanced accuracy, and
0.8514 macro F1 (Table~\ref{tab:comparison}). While slightly below the CNN and
Transformer on aggregate metrics, the GNN shows markedly improved handling of
the L2/3~IT class: F1 of 0.247 versus 0.011 for the CNN and near-zero for the
Transformer. This improvement is attributable to the neighborhood aggregation
mechanism---L2/3~IT cells connected to correctly-classified neighbors receive
corrective signal that overrides the raw feature-space ambiguity with Micro-PVM.

\subsection{Inductive SMART-seq Transfer}
For inductive transfer, SMART-seq node features were stacked with training
features, new cells were connected to their $k$ nearest training neighbors via
a second $k$-NN pass, and a single forward pass produced embeddings for all
nodes. The GNN achieves 88.56\% accuracy on 16,508 SMART-seq cells, substantially
higher than either the CNN or the Transformer. The graph-based connectivity
allows SMART-seq cells to borrow strength from their 10x training neighbors,
effectively providing implicit batch correction through feature-space proximity.

% ══════════════════════════════════════════════════════════════════════════════
\section{Cross-Model Comparison}

Table~\ref{tab:comparison} summarizes all three models on the held-out test set.

\begin{table}[H]
\centering
\caption{Test set performance comparison across all three architectures.}
\label{tab:comparison}
\begin{tabular}{@{}lrrrrc@{}}
\toprule
Model & Acc & Bal.\ Acc & Macro F1 & L2/3~IT F1 & SMART-seq Transfer \\
\midrule
1D CNN               & 0.6614 & 0.9251 & 0.8705 & 0.011          & ---                \\
Transformer (TOSICA) & 0.6656 & 0.9265 & 0.8672 & $\approx$0     & 4.8\% classified   \\
GraphSAGE GNN        & 0.6317 & 0.9005 & 0.8514 & \textbf{0.247} & \textbf{88.56\%}   \\
\bottomrule
\end{tabular}
\end{table}

\noindent The three models occupy different points on the accuracy--interpretability--transfer
tradeoff. The CNN is the simplest and achieves the highest macro F1, but provides
no interpretability and no principled mechanism for unknown-cell detection or
cross-technology transfer. The Transformer sacrifices a negligible amount of F1
for biologically grounded pathway-level interpretability and a probability-based
rejection option, but the rejection rate on SMART-seq is too aggressive to be
practically useful without explicit batch correction. The GNN partially resolves
the L2/3~IT failure through neighborhood context and achieves the strongest
inductive transfer by leveraging feature-space proximity, but its full-batch
training regime limits scalability to larger datasets.

% ══════════════════════════════════════════════════════════════════════════════
\section{Shortcomings and Discussion}

The results across all three architectures expose a consistent theme: the
L2/3~IT/Micro-PVM misclassification is a data-level pathology, not a modeling
failure. The L2/3~IT class is simultaneously the largest (31.7\% of cells) and,
in the HVG feature space, the most ambiguous---its expression profile overlaps
substantially with Micro-PVM after variance-based gene selection. Class-weighted
loss, focal loss, and neighborhood aggregation all partially mitigate but do not
eliminate this overlap. A coarse-to-fine hierarchical classifier, or a dedicated
binary disambiguator trained specifically on L2/3~IT vs.\ Micro-PVM, is likely
necessary to fully resolve this failure mode.

The Transformer's high SMART-seq rejection rate (95.2\%) highlights a tension
between the interpretability benefits of pathway tokenization and batch
robustness: while pathway masks suppress individual gene noise, the overall
expression distribution shift between 10x and SMART-seq is large enough to push
most cells below the 0.95 confidence threshold. The GNN's strong transfer
performance (88.56\%) suggests that explicit structural connectivity---anchoring
new cells to their nearest reference neighbors---is more effective at bridging
technology gaps than representation-level interventions alone.

% ══════════════════════════════════════════════════════════════════════════════
\section{Conclusion}

We evaluated three deep learning architectures for single-cell RNA-seq cell type
classification on the Allen Brain Atlas Human M1 dataset. A 1D CNN achieves
66.14\% accuracy (92.51\% balanced) and is the strongest on per-class F1 for
the majority of cell types, but catastrophically fails on the dominant L2/3~IT
class. A pathway-masked Transformer marginally improves aggregate accuracy to
66.56\% while adding biologically interpretable attention scores and a
probability-based unknown-cell filter; its UMAP on attention weights reveals
pathway-level structure invisible to PCA on raw gene expression. A GraphSAGE
GNN achieves 63.17\% accuracy but partially resolves the L2/3~IT failure
(F1 0.247) and achieves the strongest cross-technology transfer (88.56\% on
16,508 SMART-seq cells) through inductive neighborhood connectivity. The
persistent L2/3~IT/Micro-PVM confusion across all three models confirms that
this failure is intrinsic to the feature representation and class distribution,
pointing to hierarchical classification or targeted data augmentation as the
most promising directions for future work. Stream processing via Polars proved
essential for handling the 30\,GB expression matrix under memory constraints,
and the truncated-download incident highlights the importance of data validation
before model evaluation.

\end{document}
